\section{Constraint Engine}
\label{sec:constraint-engine}

The constraint engine determines whether an execution plan preserves the
semantic correctness of the original tensor computation. Each candidate
passes through a sequential chain of constraint checks with early-exit:
the first violation short-circuits the chain.

\subsection{The Seven Constraint Types}

\paragraph{1. Shape.}
For each operation $o$ with inputs $t_1, \ldots, t_n$ and output
$t_{\text{out}}$, the shape-compatibility predicate
$\phi_o(\text{shape}(t_1), \ldots, \text{shape}(t_n)) \to
\text{shape}(t_{\text{out}})$ must hold (e.g.\ inner dimensions of a
matrix product must agree). Static/dynamic check, $O(|V|)$ per plan,
backend-independent.

\paragraph{2. Tensor layout.}
The layout assignment $\Lambda: \mathcal{T} \to \{\text{NC}, \text{CN},
\text{NCHW}, \text{NHWC}, \ldots\}$ must be compatible with every
consuming operation and the target backend; a kernel requiring NHWC
invalidates a plan assigning NCHW without an interposed transpose.
Static, $O(|V|)$, backend-dependent.

\paragraph{3. Memory.}
$C_{\text{mem}}(P) = 1$ iff $M(P) \leq M_{\max}(\beta)$, where $M(P)$
is the peak of the sum of live allocations over the execution schedule.
Static, $O(|V|)$, user-configurable budget.

\paragraph{4. Backend capability.}
The capability predicate $\kappa: \mathcal{O} \times \mathcal{B} \to
\{0,1\}$ must accept every operation, covering both operation-level and
datatype-level support. Static, $O(|V|)$, backend-dependent.

\paragraph{5. Numerical error.}
The estimated cumulative error must satisfy
$\|\hat{\mathbf{y}} - \mathbf{y}_{\text{ref}}\|_2 \leq \epsilon$
(default $\epsilon = 10^{-4}$). Because the reference output is
unavailable at validation time, \gate{} bounds the effect of dtype
conversions, reorderings, and layout transformations with an analytical
interval-arithmetic model; in the empirical reference implementations
(Section~\ref{subsec:empirical}) the estimate is grounded by one
calibration execution per candidate. $O(|V| \cdot d)$, tolerance
user-configurable.

\paragraph{6. Determinism.}
Rejects plans containing non-deterministic operations when
\texttt{require\_determinism} is set: atomic parallel reductions with
unstable summation order, floating-point reassociation, and inherently
non-deterministic algorithms (top-$k$ tie-breaking, sampling). Static,
$O(|V|)$.

\paragraph{7. Safety.}
Verifies that (1) every tensor read is within allocated bounds, (2)
every write targets a valid buffer, (3) buffer lifetimes cover all
uses, and (4) no aliasing creates data races. $O(|V|^2)$ worst case
(alias and liveness analysis), $O(|V|)$ in practice.

\subsection{Composability}

Adding a constraint $C_{k+1}$ yields the validator
$V'(P) = V(P) \cdot C_{k+1}(P)$ with no modification of existing
constraints---open for extension, closed for modification. Each
constraint is an independent object behind a uniform interface
(\texttt{Constraint} trait in Rust, protocol in Python), and
Theorem~\ref{thm:soundness} holds for any $k$.

\subsection{Default Validation Order}

Constraints are ordered cheapest-and-most-rejective first to maximize
early-exit benefit: shape, tensor layout, memory, backend capability,
numerical error, determinism, safety. Empirically this ordering reduces
effective validation cost by $2.5$--$3.5\times$
(Section~\ref{sec:complexity}).
